\section*{Hoofdstuk \ref{ch:b3:petrological-thermodynamics} --- Thermodynamica van gesteenten}

\begin{solution}{exo:b3:petrological-thermodynamics:1}
$\Delta_rH = -3029.3 - 910.7 + 3935.0 = \qty{-5.0}{kJ}$; $\Delta_rS = 133.5 +
41.5 - 207.4 = \qty{-32.4}{J.K^{-1}}$; $\Delta_rV = 60.40 + 22.69 - 100.07 =
\qty{-16.98}{cm^3} = \qty{-1.698}{J.bar^{-1}}$. Helling $32.4/1.698 =
\qty{19.1}{bar/K}$. Bij \qty{873}{K}: $P = 1 + (873.15\times(-32.4) +
5000)/(-1.698) = \qty{13.7}{kbar}$.
\end{solution}

\begin{solution}{exo:b3:petrological-thermodynamics:2}
$c = 2$ en $p = 3 > c$: de regel wordt geschonden. De drie fasen coëxisteren alleen
op de univariante lijn van \ce{Mg2SiO4 + SiO2 <=> 2 MgSiO3}; een willekeurig
gekozen gesteente ligt daar niet, dus de assemblage registreert een onvolledige
reactie, bijvoorbeeld forsteriet dat omhuld is door een rand van enstatiet en
zo van het kwarts is afgeschermd.
\end{solution}

\begin{solution}{exo:b3:petrological-thermodynamics:3}
$a_{\text{grossulaar}} = 0.20^3 = 0.008$; $a_{\text{pyroop}} = 0.15^3 =
\num{3.4e-3}$; $a_{\text{anorthiet}} = 0.40$.
\end{solution}

\begin{solution}{exo:b3:petrological-thermodynamics:4}
$T_c = 15\,000/(2\times8.314) = \qty{902}{K}$; $W = 2\times8.314\times1000
= \qty{16.6}{kJ.mol^{-1}}$.
\end{solution}

\begin{solution}{exo:b3:petrological-thermodynamics:5}
$\Delta_rG = -5432.2 + 2\times(-2088.8) - 730.2 + 3\times3478.1 =
\qty{94.3}{kJ}$; $P = 1 + 94\,300/6.61 = \qty{14.3}{kbar}$. De experimentele
lijn geeft $22.80\times726.85 - 1093 = \qty{15.5}{kbar}$, \qty{1.2}{kbar}
hoger.
\end{solution}

\begin{solution}{exo:b3:petrological-thermodynamics:6}
$RT = \qty{7483}{J.mol^{-1}}$. Bij $x = 0.5$: $7483\ln0.5 + 20\,000/4 = -5187 +
5000 = \qty{-187}{J.mol^{-1}}$. Bij $x = 0.111$: $7483\times(0.111\ln0.111 +
0.889\ln0.889) + 20\,000\times0.111\times0.889 = -2609 + 1974 =
\qty{-635}{J.mol^{-1}}$, en hetzelfde bij 0.889 door symmetrie. Gelijke
hoeveelheden van de twee flanken hebben de bulksamenstelling 0.5 en de Gibbs-energie
\qty{-635}{J.mol^{-1}}: splitsen verlaagt de energie met \qty{448}{J.mol^{-1}}.
\end{solution}

\begin{solution}{exo:b3:petrological-thermodynamics:7}
$c + 2 = 5$ fasen en 5 lijnen. (D) $A + C = 2B$: $0 + 2/3 = 2\times1/3$. (A)
$B + D = 2C$: $1/3 + 1 = 2\times2/3$. (C) $2A + D = 3B$: $0 + 1 = 3\times1/3$.
(B) $A + 2D = 3C$: $0 + 2 = 3\times2/3$.
\end{solution}

\begin{solution}{exo:b3:petrological-thermodynamics:8}
Volgens de hefboomregel maakt de fase bij 0.889 $(0.3 - 0.111)/(0.889 - 0.111) =
0.243$, \qty{24}{\percent}, en de fase bij 0.111 \qty{76}{\percent}.
\end{solution}

\begin{solution}{exo:b3:petrological-thermodynamics:9}
$\partial P/\partial X_{\text{an}} = -3RT/(6.21X_{\text{an}}) =
-3\times8.314\times900/(6.21\times0.30) = \qty{-12000}{bar}$ per eenheid van
$X_{\text{an}}$: een fout van 0.01 verschuift de druk met \qty{120}{bar}.
\end{solution}

\begin{solution}{exo:b3:petrological-thermodynamics:10}
$G'_{\text{mix}} = RT\ln[x/(1 - x)] + W(1 - 2x)$ en $G''_{\text{mix}} =
RT/[x(1 - x)] - 2W$, nul waar $x(1 - x) = RT/(2W)$. Voor $x = 0.3$: $T =
2\times20\,000\times0.21/8.314 = \qty{1010}{K}$; het \omterm{def:b3:petrological-thermodynamics:solvus}{solvus} wordt bereikt bij
$20\,000\times0.4/(8.314\ln(7/3)) = \qty{1136}{K}$. Tussen beide in is het
kristal metastabiel: splitsen zou de energie verlagen, maar een kleine
fluctuatie verhoogt haar, dus moet de tweede fase nucleëren.
\end{solution}

\begin{solution}{exo:b3:petrological-thermodynamics:11}
Omdat $G_{\text{mix}}(x) = G_{\text{mix}}(1 - x)$, raakt de gemeenschappelijke
raaklijn de kromme in twee symmetrische punten $x$ en $1 - x$ en is ze horizontaal:
$G'_{\text{mix}}(x) = 0$, wat de vergelijking is. Oplossen geeft $T(x) = W(1 -
2x)/[R\ln((1 - x)/x)]$. Met $x = 1/2 - \varepsilon$ geldt $1 - 2x = 2\varepsilon$
en $\ln((1/2 + \varepsilon)/(1/2 - \varepsilon)) \approx 4\varepsilon$, dus
$T \to W/(2R)$ voor $\varepsilon \to 0$; en voor $x < 1/2$ neemt $T(x)$ toe naarmate $x$
nadert tot 1/2, dus deze grenswaarde is het maximum.
\end{solution}

\begin{solution}{exo:b3:petrological-thermodynamics:12}
$(2.0 - 0.08\times10)/0.92 = \qty{1.30}{wt\percent}$. Mangaan begunstigt
granaat: met minder ervan in de \omterm{def:b3:petrological-thermodynamics:pseudosection}{effectieve bulksamenstelling} groeien de randen in
een ander veld en langs andere isopleeten dan de totaalanalyse van het gesteente
zou voorspellen.
\end{solution}

\begin{solution}{pb:b3:petrological-thermodynamics:1}
\textbf{1.} $\Delta_rH^\circ = -6640.0 + 2\times(-2593.8) - 910.7 + 3\times4234.0
= \qty{-36.3}{kJ}$; $\Delta_rS^\circ = 260.1 + 165.6 + 41.5 - 597.9 =
\qty{-130.7}{J.K^{-1}}$; $\Delta_rV = 125.28 + 88.30 + 22.69 - 302.37 =
\qty{-66.10}{cm^3} = \qty{-6.61}{J.bar^{-1}}$.

\textbf{2.} $-36\,300 + 298.15\times130.7 = \qty{+2.67}{kJ}$: positief, dus
anorthiet is stabiel aan het oppervlak.

\textbf{3.} $P = 1 + (900\times(-130.7) + 36\,300)/(-6.61) = \qty{12.3}{kbar}$;
helling $130.7/6.61 = \qty{19.8}{bar/K}$.

\textbf{4.} $1 + 81\,500/6.61 = \qty{12.3}{kbar}$ bij \qty{900}{K}; $1 +
119\,800/6.61 = \qty{18.1}{kbar}$ bij \qty{1200}{K}.

\textbf{5.} Bij 626.85 en \qty{926.85}{\celsius} geeft de experimentele lijn
13.2 en \qty{20.0}{kbar}: 0.9 en \qty{1.9}{kbar} boven de tabellen.

{\sloppy\textbf{6.} $1914\times6.61 = \qty{12.7}{kJ}$. De gecombineerde onzekerheid is
$\sqrt{3.2^2 + (3\times2.0)^2 + (2\times2.0)^2 + 1.0^2} = \qty{8.0}{kJ}$: het
verschil is 1.6 keer dat, mogelijk voor calorimetrie, en de reden waarom de lijn
uit de experimenten wordt genomen.\par}

\textbf{7.} $\Delta_rC_p$ is klein bij vaste stoffen, en de twee integralen
heffen elkaar gedeeltelijk op: de tabellen verschuiven de lijn met slechts 26 en \qty{112}{bar}.

\textbf{8.} \ce{SiO2}, \ce{Al2O3}, \ce{FeO}, \ce{MgO}, \ce{MnO}, \ce{CaO},
\ce{Na2O}, \ce{K2O} en \ce{H2O}: $c = 9$. Fasen: granaat, biotiet,
muscoviet, plagioklaas, kyaniet, kwarts en fluïdum, $p = 7$.

\textbf{9.} $F = 9 - 7 + 2 = 4$: druk, temperatuur en twee
samenstellingsvariabelen. $p = 7 \le c = 9$: de regel geldt, en de associatie
kan over een gebied van het druk--temperatuurvlak stabiel zijn.

\textbf{10.} Twee zuivere \ce{Al2SiO5}-polymorfen in evenwicht liggen op de
univariante lijn kyaniet = sillimaniet: één betrekking tussen $P$ en $T$, en
een thermometer levert dan de druk.

\textbf{11.} Het fluïdum legt de chemische potentiaal van water op: 8 vrije
componenten voor 6 vaste fasen, $F = 8 - 6 + 2 = 4$, en $p \le c$
geldt nog steeds.

\textbf{12.} $a_{\text{grossulaar}} = 0.10^3 = 0.001$, $a_{\text{anorthiet}} =
0.30$; $\ln K = \ln(0.001/0.30^3) = -3.30$.

\textbf{13.} $K_D = (0.12/0.74)/(0.45/0.55) = 0.1622/0.8182 = 0.198$.

\textbf{14.} $T = (52\,635 + 2488)/(20.20 + 24.94\times1.618) = 55\,123/60.57
= \qty{910}{K}$, of \qty{637}{\celsius}.

\textbf{15.} $P_0 = 22.80\times636.9 - 1093 = \qty{13429}{bar}$; $(3\times8.314
\times910.1/6.21)\ln(1/3) = \qty{-4016}{bar}$; $P = \qty{9.41}{kbar}$.

\textbf{16.} Bij \qty{9.41}{kbar}, \qty{644.2}{\celsius}, daarna
\qty{9.55}{kbar}; bij die druk \qty{644.9}{\celsius} en
\qty{9.56}{kbar}, daarna ongewijzigd: \qty{645}{\celsius} en
\qty{9.56}{kbar}.

\textbf{17.} \qty{641}{\celsius} en \qty{8.79}{kbar}. De tabellen leggen de
eindlid-lijn ongeveer \qty{1}{kbar} te laag (Deel I); de experimentele lijn
wordt behouden.

\textbf{18.} Met $X_{\text{an}} = 0.35$: \qty{642}{\celsius} en
\qty{8.93}{kbar}, \qty{0.63}{kbar} minder.

\textbf{19.} Bij \qty{918}{K}: $P = 1 + (918\times12.6 - 7700)/0.571 =
\qty{6.8}{kbar}$, onder \qty{9.56}{kbar}: kyaniet, zoals waargenomen.

\textbf{20.} $z = \num{9.56e8}/(2800\times9.81) = \qty{34.8}{km}$.

\textbf{21.} $\Delta P = (3\times8.314\times918/6.21)\ln(0.07/0.10) =
\qty{-1.3}{kbar}$: de kern groeide ongeveer \qty{4.8}{km} ondieper. De granaat
groeide tijdens begraving, en zijn rand registreert het diepste punt.

\textbf{22.} De Varden-schist bereikte een piekdruk van
\textbf{\qty[text-series-to-math=true]{9.6}{kbar}} bij
\qty{645}{\celsius}, een begraving van ongeveer \qty{35}{km}.
\end{solution}
